% Register file virtualisation (Lesson 6).  Six architectural register
% names (R1 of threads 1-3, R2 of threads 4-6) whose values are live in
% disjoint intervals share two physical registers: P0 and P1.
% Redraws the outline figure m6l2_opt2 (lifetime bars added).
% Usage: % Register file virtualisation (Lesson 6).  Six architectural register
% names (R1 of threads 1-3, R2 of threads 4-6) whose values are live in
% disjoint intervals share two physical registers: P0 and P1.
% Redraws the outline figure m6l2_opt2 (lifetime bars added).
% Usage: % Register file virtualisation (Lesson 6).  Six architectural register
% names (R1 of threads 1-3, R2 of threads 4-6) whose values are live in
% disjoint intervals share two physical registers: P0 and P1.
% Redraws the outline figure m6l2_opt2 (lifetime bars added).
% Usage: % Register file virtualisation (Lesson 6).  Six architectural register
% names (R1 of threads 1-3, R2 of threads 4-6) whose values are live in
% disjoint intervals share two physical registers: P0 and P1.
% Redraws the outline figure m6l2_opt2 (lifetime bars added).
% Usage: \input{figures/l06-rf-virtualization}
% Width: body (120 mm); the drawing is about 112 mm wide.
\begin{tikzpicture}[x=1mm, y=1mm, font=\scriptsize\sffamily,
    >={Stealth[length=1.6mm]}]
  \node[anchor=south west, font=\scriptsize\sffamily\bfseries] at (0,40)
    {\iflangja{名前空間}{name space}};
  \node[anchor=south west, font=\scriptsize\sffamily\bfseries] at (18,40)
    {\iflangja{値が生存する区間}{interval in which the value is live}};
  \node[anchor=south west, font=\scriptsize\sffamily\bfseries] at (90,40)
    {\iflangja{物理レジスタ}{physical registers}};
  \draw[ln-rule] (18,2) -- (82,2);
  \draw[->, ln-muted] (18,0) -- (82,0) node[right] {\iflangja{時間}{time}};
  \foreach \n/\a/\b/\c [count=\i from 0] in
      {Th1.R1/0/16/ln-accent!35, Th2.R1/20/36/ln-accent!35, Th3.R1/42/60/ln-accent!35,
       Th4.R2/6/22/ln-tip!35, Th5.R2/26/42/ln-tip!35, Th6.R2/46/62/ln-tip!35} {
    \node[anchor=west, font=\scriptsize\ttfamily] at (0,36-\i*6) {\n};
    \draw[ln-rule, dotted] (18,36-\i*6) -- (82,36-\i*6);
    \draw[fill=\c] (18+\a,34.7-\i*6) rectangle (18+\b,37.3-\i*6);
  }
  \draw[fill=ln-accent!35] (92,27) rectangle ++(14,6) node[midway] {P0};
  \draw[fill=ln-tip!35] (92,9) rectangle ++(14,6) node[midway] {P1};
  \draw[ln-muted] (83,38) -- ++(1.5,0) -- ++(0,-14) -- ++(-1.5,0);
  \draw[ln-muted] (83,20) -- ++(1.5,0) -- ++(0,-14) -- ++(-1.5,0);
  \draw[->] (84.5,31) -- (92,30);
  \draw[->] (84.5,13) -- (92,12);
\end{tikzpicture}

% Width: body (120 mm); the drawing is about 112 mm wide.
\begin{tikzpicture}[x=1mm, y=1mm, font=\scriptsize\sffamily,
    >={Stealth[length=1.6mm]}]
  \node[anchor=south west, font=\scriptsize\sffamily\bfseries] at (0,40)
    {\iflangja{名前空間}{name space}};
  \node[anchor=south west, font=\scriptsize\sffamily\bfseries] at (18,40)
    {\iflangja{値が生存する区間}{interval in which the value is live}};
  \node[anchor=south west, font=\scriptsize\sffamily\bfseries] at (90,40)
    {\iflangja{物理レジスタ}{physical registers}};
  \draw[ln-rule] (18,2) -- (82,2);
  \draw[->, ln-muted] (18,0) -- (82,0) node[right] {\iflangja{時間}{time}};
  \foreach \n/\a/\b/\c [count=\i from 0] in
      {Th1.R1/0/16/ln-accent!35, Th2.R1/20/36/ln-accent!35, Th3.R1/42/60/ln-accent!35,
       Th4.R2/6/22/ln-tip!35, Th5.R2/26/42/ln-tip!35, Th6.R2/46/62/ln-tip!35} {
    \node[anchor=west, font=\scriptsize\ttfamily] at (0,36-\i*6) {\n};
    \draw[ln-rule, dotted] (18,36-\i*6) -- (82,36-\i*6);
    \draw[fill=\c] (18+\a,34.7-\i*6) rectangle (18+\b,37.3-\i*6);
  }
  \draw[fill=ln-accent!35] (92,27) rectangle ++(14,6) node[midway] {P0};
  \draw[fill=ln-tip!35] (92,9) rectangle ++(14,6) node[midway] {P1};
  \draw[ln-muted] (83,38) -- ++(1.5,0) -- ++(0,-14) -- ++(-1.5,0);
  \draw[ln-muted] (83,20) -- ++(1.5,0) -- ++(0,-14) -- ++(-1.5,0);
  \draw[->] (84.5,31) -- (92,30);
  \draw[->] (84.5,13) -- (92,12);
\end{tikzpicture}

% Width: body (120 mm); the drawing is about 112 mm wide.
\begin{tikzpicture}[x=1mm, y=1mm, font=\scriptsize\sffamily,
    >={Stealth[length=1.6mm]}]
  \node[anchor=south west, font=\scriptsize\sffamily\bfseries] at (0,40)
    {\iflangja{名前空間}{name space}};
  \node[anchor=south west, font=\scriptsize\sffamily\bfseries] at (18,40)
    {\iflangja{値が生存する区間}{interval in which the value is live}};
  \node[anchor=south west, font=\scriptsize\sffamily\bfseries] at (90,40)
    {\iflangja{物理レジスタ}{physical registers}};
  \draw[ln-rule] (18,2) -- (82,2);
  \draw[->, ln-muted] (18,0) -- (82,0) node[right] {\iflangja{時間}{time}};
  \foreach \n/\a/\b/\c [count=\i from 0] in
      {Th1.R1/0/16/ln-accent!35, Th2.R1/20/36/ln-accent!35, Th3.R1/42/60/ln-accent!35,
       Th4.R2/6/22/ln-tip!35, Th5.R2/26/42/ln-tip!35, Th6.R2/46/62/ln-tip!35} {
    \node[anchor=west, font=\scriptsize\ttfamily] at (0,36-\i*6) {\n};
    \draw[ln-rule, dotted] (18,36-\i*6) -- (82,36-\i*6);
    \draw[fill=\c] (18+\a,34.7-\i*6) rectangle (18+\b,37.3-\i*6);
  }
  \draw[fill=ln-accent!35] (92,27) rectangle ++(14,6) node[midway] {P0};
  \draw[fill=ln-tip!35] (92,9) rectangle ++(14,6) node[midway] {P1};
  \draw[ln-muted] (83,38) -- ++(1.5,0) -- ++(0,-14) -- ++(-1.5,0);
  \draw[ln-muted] (83,20) -- ++(1.5,0) -- ++(0,-14) -- ++(-1.5,0);
  \draw[->] (84.5,31) -- (92,30);
  \draw[->] (84.5,13) -- (92,12);
\end{tikzpicture}

% Width: body (120 mm); the drawing is about 112 mm wide.
\begin{tikzpicture}[x=1mm, y=1mm, font=\scriptsize\sffamily,
    >={Stealth[length=1.6mm]}]
  \node[anchor=south west, font=\scriptsize\sffamily\bfseries] at (0,40)
    {\iflangja{名前空間}{name space}};
  \node[anchor=south west, font=\scriptsize\sffamily\bfseries] at (18,40)
    {\iflangja{値が生存する区間}{interval in which the value is live}};
  \node[anchor=south west, font=\scriptsize\sffamily\bfseries] at (90,40)
    {\iflangja{物理レジスタ}{physical registers}};
  \draw[ln-rule] (18,2) -- (82,2);
  \draw[->, ln-muted] (18,0) -- (82,0) node[right] {\iflangja{時間}{time}};
  \foreach \n/\a/\b/\c [count=\i from 0] in
      {Th1.R1/0/16/ln-accent!35, Th2.R1/20/36/ln-accent!35, Th3.R1/42/60/ln-accent!35,
       Th4.R2/6/22/ln-tip!35, Th5.R2/26/42/ln-tip!35, Th6.R2/46/62/ln-tip!35} {
    \node[anchor=west, font=\scriptsize\ttfamily] at (0,36-\i*6) {\n};
    \draw[ln-rule, dotted] (18,36-\i*6) -- (82,36-\i*6);
    \draw[fill=\c] (18+\a,34.7-\i*6) rectangle (18+\b,37.3-\i*6);
  }
  \draw[fill=ln-accent!35] (92,27) rectangle ++(14,6) node[midway] {P0};
  \draw[fill=ln-tip!35] (92,9) rectangle ++(14,6) node[midway] {P1};
  \draw[ln-muted] (83,38) -- ++(1.5,0) -- ++(0,-14) -- ++(-1.5,0);
  \draw[ln-muted] (83,20) -- ++(1.5,0) -- ++(0,-14) -- ++(-1.5,0);
  \draw[->] (84.5,31) -- (92,30);
  \draw[->] (84.5,13) -- (92,12);
\end{tikzpicture}
